% !TEX program = xelatex
\documentclass[11pt,a4paper]{article}

\usepackage[a4paper,margin=2.2cm]{geometry}
\usepackage{fontspec}
\setmainfont{TeX Gyre Termes}
\setmonofont{DejaVu Sans Mono}
\usepackage{amsmath,amssymb}
\usepackage{booktabs,tabularx,array}
\usepackage{enumitem}
\usepackage{xcolor}
\usepackage{hyperref}
\usepackage{titlesec}

\definecolor{ink}{HTML}{1F2937}
\hypersetup{colorlinks=true, linkcolor=ink, urlcolor=ink}
\setlength{\parindent}{0pt}
\setlength{\parskip}{6pt}
\renewcommand{\arraystretch}{1.15}
\titleformat{\section}{\Large\bfseries\sffamily\color{ink}}{\thesection}{0.7em}{}
\titleformat{\subsection}{\large\bfseries\sffamily\color{ink}}{\thesubsection}{0.7em}{}

\newcommand{\code}[1]{\texttt{#1}}
\newcommand{\ind}[1]{\mathbf{1}\{#1\}}

\title{\bfseries\sffamily Công thức trích xuất đặc trưng cho RawML trong Online Caching}
\author{Dự án \code{tin03092006-cell/Online-Caching}}
\date{\today}

\begin{document}
\maketitle

\section{Mục đích}

Tài liệu này bổ sung công thức toán học cho bước trích xuất đặc trưng trong pipeline online caching. Các đặc trưng này được dùng để huấn luyện expert học máy \code{RawML}, cụ thể là mô hình hồi quy dự đoán khoảng cách đến lần truy cập tiếp theo của một item trong cache.

Trong code, các feature chính gồm:
\[
    \code{recency},\quad
    \code{frequency},\quad
    \code{recent\_frequency},\quad
    \code{average\_inter\_arrival},\quad
    \code{cache\_age}.
\]
Biến mục tiêu là:
\[
    \code{target\_next\_distance}.
\]

\section{Ký hiệu chung}

Cho trace request:
\[
    \mathcal{R} = (r_0,r_1,\ldots,r_{T-1}),
\]
trong đó \(r_t\) là item được request tại thời điểm \(t\). Ta dùng chỉ số bắt đầu từ 0 để khớp với code Python.

Gọi:
\[
    C_t
\]
là tập item đang nằm trong cache ngay trước khi xử lý request \(r_t\).

Gọi \(W\) là kích thước cửa sổ gần đây, tương ứng với tham số:
\[
    W = \code{recent\_window\_size}.
\]

Khi cần tạo một dòng dữ liệu huấn luyện tại thời điểm \(t\), pipeline xét từng item:
\[
    x \in C_t.
\]
Với mỗi item \(x\), ta tạo vector feature:
\[
    \phi_t(x)
    =
    \left[
    R_t(x),
    F_t(x),
    RF_t(x),
    IA_t(x),
    A_t(x)
    \right].
\]

\section{Feature 1: Recency}

\subsection{Định nghĩa}

Gọi lần truy cập gần nhất của item \(x\) trước thời điểm \(t\) là:
\[
    L_t(x)
    =
    \max\{\tau < t : r_{\tau}=x\}.
\]

Nếu item \(x\) đã từng xuất hiện trước thời điểm \(t\), recency là:
\[
    R_t(x) = t - L_t(x).
\]

Nếu item \(x\) chưa từng xuất hiện, code dùng giá trị:
\[
    R_t(x) = t + 1.
\]

Vậy công thức đầy đủ:
\[
    R_t(x)
    =
    \begin{cases}
        t - L_t(x), & \text{nếu tồn tại } L_t(x),\\
        t + 1, & \text{nếu } x \text{ chưa từng xuất hiện trước } t.
    \end{cases}
\]

\subsection{Ý nghĩa}

Recency đo item đã vắng mặt bao lâu. Nếu \(R_t(x)\) lớn, item \(x\) đã lâu chưa được request. Đây là tín hiệu tương ứng với trực giác của LRU.

\section{Feature 2: Frequency}

\subsection{Định nghĩa}

Frequency toàn cục của item \(x\) trước thời điểm \(t\) là:
\[
    F_t(x)
    =
    \sum_{\tau=0}^{t-1}
    \ind{r_{\tau}=x}.
\]

Trong code, đại lượng này tương ứng với:
\[
    \code{access\_counts[x]}.
\]

\subsection{Ý nghĩa}

Frequency đo item đã được request bao nhiêu lần trong quá khứ. Nếu \(F_t(x)\) lớn, item có thể là item phổ biến. Đây là tín hiệu gần với trực giác của LFU, nhưng cần phân biệt:

\begin{itemize}[leftmargin=1.4em]
    \item \code{frequency} trong feature là tần suất toàn cục.
    \item LFU expert trong cache thường dùng số lần truy cập kể từ khi item được đưa vào cache hiện tại.
\end{itemize}

\section{Feature 3: Recent frequency}

\subsection{Định nghĩa}

Recent frequency đếm số lần item \(x\) xuất hiện trong cửa sổ gần đây có độ dài \(W\). Cửa sổ trước thời điểm \(t\) là:
\[
    [\max(0,t-W),\,t-1].
\]

Do đó:
\[
    RF_t(x)
    =
    \sum_{\tau=\max(0,t-W)}^{t-1}
    \ind{r_{\tau}=x}.
\]

\subsection{Ý nghĩa}

Recent frequency khác frequency toàn cục ở chỗ nó chỉ nhìn vào đoạn gần đây. Nó giúp mô hình phát hiện item đang hot trong thời gian ngắn, kể cả khi tổng frequency trong toàn bộ quá khứ chưa lớn.

Nếu \(F_t(x)\) lớn nhưng \(RF_t(x)\) nhỏ, item có thể là item từng hot nhưng hiện đã nguội. Nếu \(RF_t(x)\) lớn, item có thể đang nằm trong một pha truy cập mạnh.

\section{Feature 4: Average inter-arrival time}

\subsection{Định nghĩa}

Gọi các thời điểm item \(x\) xuất hiện trước \(t\) là:
\[
    p_1 < p_2 < \cdots < p_m < t.
\]

Các khoảng cách giữa hai lần xuất hiện liên tiếp là:
\[
    g_i = p_i - p_{i-1}, \qquad i=2,\ldots,m.
\]

Nếu \(m\ge 2\), average inter-arrival time là:
\[
    IA_t(x)
    =
    \frac{1}{m-1}
    \sum_{i=2}^{m}
    (p_i-p_{i-1}).
\]

Nếu \(m<2\), tức chưa đủ hai lần truy cập để tính khoảng cách trung bình, code dùng recency làm giá trị thay thế:
\[
    IA_t(x) = R_t(x).
\]

Công thức đầy đủ:
\[
    IA_t(x)
    =
    \begin{cases}
        \dfrac{1}{m-1}\displaystyle\sum_{i=2}^{m}(p_i-p_{i-1}), & m\ge 2,\\[1.2em]
        R_t(x), & m<2.
    \end{cases}
\]

\subsection{Ý nghĩa}

Average inter-arrival time mô tả chu kỳ truy cập của item. Nếu \(IA_t(x)\) nhỏ, item có xu hướng quay lại thường xuyên. Nếu \(IA_t(x)\) lớn, item thường mất lâu mới quay lại.

Đặc trưng này giúp mô hình phân biệt hai item có cùng frequency nhưng chu kỳ xuất hiện khác nhau.

\section{Feature 5: Cache age}

\subsection{Định nghĩa}

Gọi thời điểm item \(x\) được đưa vào cache hiện tại là:
\[
    I_t(x).
\]

Cache age là:
\[
    A_t(x)
    =
    t - I_t(x).
\]

Nếu vì lý do kỹ thuật không tìm thấy thời điểm insert, code dùng mặc định \(I_t(x)=t\), khi đó:
\[
    A_t(x)=0.
\]

\subsection{Ý nghĩa}

Cache age đo item đã chiếm vị trí trong cache bao lâu. Nếu \(A_t(x)\) lớn nhưng item ít được truy cập, item đó có thể đang giữ cache một cách không hiệu quả.

\section{Biến phụ trợ: cache access count cho LFU}

LFU expert không nhất thiết dùng đúng feature \code{frequency}. Trong implementation, LFU dựa trên số lần item được truy cập kể từ khi nó được đưa vào cache hiện tại.

Gọi:
\[
    I_t(x)
\]
là thời điểm item \(x\) được đưa vào cache hiện tại. Khi đó cache-local frequency là:
\[
    CF_t(x)
    =
    \sum_{\tau=I_t(x)}^{t-1}
    \ind{r_{\tau}=x}.
\]

LFU chọn:
\[
    e_t^{\mathrm{LFU}}
    =
    \arg\min_{x\in C_t}
    CF_t(x).
\]

Nếu hòa, implementation dùng thêm last access time và item ID để phá hòa.

\section{Target: Next access distance}

\subsection{Định nghĩa}

Biến mục tiêu của \code{RawML} là khoảng cách đến lần truy cập tiếp theo.

Gọi lần xuất hiện tiếp theo của item \(x\) sau thời điểm \(t\) là:
\[
    N_t(x)
    =
    \min\{\tau>t : r_{\tau}=x\}.
\]

Nếu tồn tại \(N_t(x)\), target là:
\[
    D_t^+(x)
    =
    N_t(x)-t.
\]

Nếu không còn lần truy cập nào trong tương lai, code gán:
\[
    D_t^+(x)=T+1.
\]

Công thức đầy đủ:
\[
    D_t^+(x)
    =
    \begin{cases}
        N_t(x)-t, & \text{nếu tồn tại } N_t(x),\\
        T+1, & \text{nếu } x \text{ không xuất hiện lại trong tương lai}.
    \end{cases}
\]

\subsection{Ý nghĩa}

Target này chính là đại lượng mà Belady/OPT biết chính xác. Nếu \(D_t^+(x)\) lớn, item \(x\) còn lâu mới được dùng lại, nên nó là ứng viên tốt để evict.

\code{RawML} cố gắng học hàm:
\[
    f_{\theta}(\phi_t(x))
    \approx
    D_t^+(x).
\]

Khi chạy cache online, \code{RawML} chọn item có dự đoán lớn nhất:
\[
    e_t^{\mathrm{RawML}}
    =
    \arg\max_{x\in C_t}
    f_{\theta}(\phi_t(x)).
\]

\section{Cách tạo một dòng training sample}

Tại thời điểm \(t\), nếu cache đã đầy, với mỗi item \(x\in C_t\), pipeline tạo một dòng dữ liệu:

\[
    \left(
    R_t(x),
    F_t(x),
    RF_t(x),
    IA_t(x),
    A_t(x),
    D_t^+(x)
    \right).
\]

Trong đó:

\[
    \underbrace{
    \left(
    R_t(x),
    F_t(x),
    RF_t(x),
    IA_t(x),
    A_t(x)
    \right)
    }_{\text{features }X}
    \quad
    \longrightarrow
    \quad
    \underbrace{D_t^+(x)}_{\text{label }y}.
\]

Bảng huấn luyện có dạng:

\begin{center}
\begin{tabular}{rrrrrr}
\toprule
recency & frequency & recent freq. & avg. inter-arrival & cache age & target \\
\midrule
\(R_t(x)\) & \(F_t(x)\) & \(RF_t(x)\) & \(IA_t(x)\) & \(A_t(x)\) & \(D_t^+(x)\) \\
\bottomrule
\end{tabular}
\end{center}

\section{Thứ tự cập nhật tại mỗi thời điểm}

Thứ tự xử lý rất quan trọng. Với request \(r_t\):

\begin{enumerate}[leftmargin=1.4em]
    \item Trước khi cập nhật lịch sử bằng \(r_t\), hệ thống có trạng thái \(F_t\).
    \item Nếu cache đã đầy và cần tạo training rows, feature của các item \(x\in C_t\) được tính từ lịch sử trước \(t\).
    \item Target \(D_t^+(x)\) được tính bằng cách nhìn vào phần tương lai của trace.
    \item Sau đó thuật toán xử lý request \(r_t\): hit hoặc miss.
    \item Cuối cùng, lịch sử được cập nhật: last access, access count, recent window, inter-arrival statistics.
\end{enumerate}

Vì vậy, feature không được dùng thông tin tương lai. Chỉ label \(D_t^+(x)\) dùng tương lai, và label này chỉ dùng khi train/evaluate mô hình, không dùng trực tiếp trong online decision.

\section{Ví dụ minh họa ngắn}

Xét trace:
\[
    \mathcal{R}=(A,B,A,C,A,B)
\]
và xét thời điểm:
\[
    t=4,\qquad r_4=A.
\]

Lịch sử trước \(t=4\) là:
\[
    (A,B,A,C).
\]

Với item \(B\):

\[
    L_4(B)=1,
    \qquad
    R_4(B)=4-1=3.
\]

Frequency:
\[
    F_4(B)=1.
\]

Nếu \(W=3\), cửa sổ gần đây là các vị trí \(1,2,3\), tức:
\[
    (B,A,C).
\]
Do đó:
\[
    RF_4(B)=1.
\]

Item \(B\) mới xuất hiện một lần trước \(t=4\), nên chưa có inter-arrival gap:
\[
    IA_4(B)=R_4(B)=3.
\]

Lần xuất hiện tiếp theo của \(B\) là tại vị trí 5:
\[
    N_4(B)=5.
\]
Do đó target:
\[
    D_4^+(B)=5-4=1.
\]

Nghĩa là tại thời điểm 4, nếu evict \(B\), chỉ một bước sau đã cần lại \(B\), nên đây là quyết định rất tệ. Một mô hình tốt nên dự đoán khoảng cách của \(B\) là nhỏ và tránh evict nó.

\section{Tóm tắt công thức}

\begin{center}
\begin{tabularx}{\textwidth}{lX}
\toprule
Đại lượng & Công thức \\
\midrule
Recency & \(R_t(x)=t-L_t(x)\), nếu chưa từng xuất hiện thì \(R_t(x)=t+1\). \\
Frequency & \(F_t(x)=\sum_{\tau=0}^{t-1}\ind{r_{\tau}=x}\). \\
Recent frequency & \(RF_t(x)=\sum_{\tau=\max(0,t-W)}^{t-1}\ind{r_{\tau}=x}\). \\
Average inter-arrival & \(IA_t(x)=\frac{1}{m-1}\sum_{i=2}^{m}(p_i-p_{i-1})\), nếu \(m<2\) thì \(IA_t(x)=R_t(x)\). \\
Cache age & \(A_t(x)=t-I_t(x)\). \\
Cache-local frequency & \(CF_t(x)=\sum_{\tau=I_t(x)}^{t-1}\ind{r_{\tau}=x}\). \\
Target next distance & \(D_t^+(x)=N_t(x)-t\), nếu không xuất hiện lại thì \(D_t^+(x)=T+1\). \\
\bottomrule
\end{tabularx}
\end{center}

\end{document}
